\documentclass{article}

\usepackage{amsmath,amssymb}
\usepackage{xurl}
\usepackage[hidelinks]{hyperref}
\setlength{\emergencystretch}{2em}

% Shared commands for notes in docs/paper-gaps/.

\DeclareUnicodeCharacter{2080}{\ifmmode{}_{0}\else\textsubscript{0}\fi}
\DeclareUnicodeCharacter{2081}{\ifmmode{}_{1}\else\textsubscript{1}\fi}
\DeclareUnicodeCharacter{2082}{\ifmmode{}_{2}\else\textsubscript{2}\fi}
\DeclareUnicodeCharacter{2099}{\ifmmode{}_{n}\else\textsubscript{n}\fi}

\newcommand{\Lean}{\textsc{Lean}}
\ExplSyntaxOn
\NewDocumentCommand{\leanid}{m}{
  \ifmmode
    \textsf{\detokenize{#1}}
  \else
    \group_begin:
    \urlstyle{sf}
    \tl_set:Nn \l_tmpa_tl {#1}
    \tl_replace_all:Nnn \l_tmpa_tl {\_} {_}
    \exp_args:NV \path \l_tmpa_tl
    \group_end:
  \fi
}
\ExplSyntaxOff
\providecommand{\tr}{\operatorname{tr}}
\newcommand{\tnleanIssueBase}{https://github.com/LionSR/TNLean/issues/}
\newcommand{\tnleanPrBase}{https://github.com/LionSR/TNLean/pull/}
\newcommand{\tnleanDocBase}{https://sirui-lu.com/TNLean/docs/}
\newcommand{\ghissue}[1]{\href{\tnleanIssueBase#1}{GitHub issue \##1}}
\newcommand{\ghpr}[1]{\href{\tnleanPrBase#1}{PR \##1}}
\newcommand{\doclink}[1]{\href{\tnleanDocBase#1.html}{\path{#1}}}

% Dirac notation and the identity operator, as in blueprint/src/macros/common.tex.
% \providecommand keeps note-local definitions authoritative.
\usepackage{dsfont}
\providecommand{\ket}[1]{|#1\rangle}
\providecommand{\bra}[1]{\langle #1|}
\providecommand{\braket}[2]{\langle #1 | #2 \rangle}
\providecommand{\ketbra}[2]{|#1\rangle\!\langle #2|}
\providecommand{\Id}{\mathds{1}}
\providecommand{\one}{\mathds{1}}
\providecommand{\C}{\mathbb{C}}
\providecommand{\MN}[1]{M_{#1}(\C)}
\providecommand{\MD}{M_D(\C)}

% Verdict marker: \gapnote{<kind>}{<status>}. Kind is the policy
% classification (clarification, local-correction, scope-restriction,
% unfaithful, false-source, open-gap); status is open, wip (elimination
% underway), resolved, or historical. Machine-read by the site index;
% prints a verdict line.
\newcommand{\gapnote}[2]{\par\noindent\textbf{Verdict:} #1 (#2).\par}


\title{Bond Dimensions of the Successive Decompositions and Minimal Resources}
\author{}
\date{2026-10-02}

\begin{document}
\maketitle
\gapnote{scope-restriction}{open}

\section{Source range}

After the successive singular value decompositions that write a state as
$\ket{\tilde\psi}=V'_{[n]}\cdots V'_{[1]}\ket{\varphi_I}$,
Ref.~\cite[lines~200--202 of the source]{Schoen2005Sequential} states:
``Simple rank considerations show that $V'_{[n-k]}$ has dimension
$2\min[D,2^k]\times\min[D,2^{k+1}]$.'' Ref.~\cite[Section~5.1,
lines~1574--1577]{PerezGarcia2007Matrix} adds that this proof ``also provides
a \emph{recipe} for the generation of any given state (with minimal
resources).'' Ref.~\cite[lines~431--458]{PerezGarcia2007Matrix} states that
every state of $N$ sites has an open-boundary representation with bond
dimension $D\le d^{\lfloor N/2\rfloor}$ satisfying three gauge conditions, and
that ``any state for which $\max_m \operatorname{rank}(\rho_m)\le D$ can be
written as a MPS of bond dimension $D$''.

\section{What is formalized}

For local dimension $d$, the isometric open-boundary representation produced
by the decompositions of a given representation has bond dimension at most
$\min(D,d^k)$ at the bond after $k$ sites counted from the end where the
decompositions start, and, after running the decompositions once from each
end, at most $\min(D_k,d^k,d^{N-k})$, where $D_k\le D$ is the bond dimension of
the given representation at that cut. Read as a statement that isometries of
the printed size suffice, the first quoted sentence is formalized with the
source's hypotheses.

For the cut $k$, let $\operatorname{rank}_k\psi$ be the rank of the matrix
$(\sigma,\tau)\mapsto\psi(\sigma\tau)$ whose rows are the configurations of
the first $k$ sites and whose columns are those of the remaining $N-k$ sites;
it is the rank of the reduced density operator $\rho_k$. The following are
formalized on chains of $N\ge 1$ sites.
\begin{enumerate}
    \item Every open-boundary representation of $\psi$ has bond dimension at
          least $\operatorname{rank}_k\psi$ at the cut $k$.
    \item If $\psi\ne0$ and $\operatorname{rank}_k\psi\le D$ for every $k$,
          successive Schmidt decompositions of $\psi$ give an open-boundary
          representation with common bound $D$ whose bond dimension at the cut
          $k$ equals $\operatorname{rank}_k\psi$ for every $k$, and in which every
          site except the last satisfies $\sum_iA_i^\dagger A_i=\mathbb 1$ (every
          site, when $\psi$ is normalized).
    \item Consequently the states generated with a $D$-dimensional ancilla,
          probabilistically or deterministically, are those whose cut ranks are
          all at most $D$ (normalized, in the deterministic case), and the least
          ancilla dimension that generates $\psi$ is $\max_k
              \operatorname{rank}_k\psi$. This is the ``minimal resources'' of the
          second quoted sentence.
    \item Every state of $N\ge 1$ sites has an open-boundary representation with
          bond dimension at most $d^{\lfloor N/2\rfloor}$, since
          $\operatorname{rank}_k\psi\le\min(d^k,d^{N-k})$.
\end{enumerate}

\section{Readings that are not formalized}

The printed size $2\min[D,2^k]\times\min[D,2^{k+1}]$ is not the minimal size
for low-rank states: a product state has $\operatorname{rank}_k\psi=1$ at every
cut, so a representation with all bond dimensions one exists, whereas
$\min[D,2^k]\ge2$ for $D\ge2$ and $k\ge1$. A representation of the printed
size still exists, since unused virtual directions can be added to a smaller
one, so the sentence is not refuted when read as asserting that isometries of
exactly that size can be chosen; it is only not minimal. Only the upper-bound
reading is formalized, and the minimal sizes are the cut ranks of the previous
section.

A single pass of the decompositions over a given representation need not reach
the cut ranks: if a bond direction of the given representation is annihilated
by the sites to its right, the left-to-right pass can keep it. The
formalization obtains the minimal bond dimensions from the Schmidt
decompositions of $\psi$ itself, as in
Ref.~\cite[lines~444--447]{PerezGarcia2007Matrix}.

The gauge conditions 1--3 of the completeness theorem
\cite[lines~431--442]{PerezGarcia2007Matrix}, namely
$\sum_iA^{[m]}_iA^{[m]\dagger}_i=\mathbb 1$, the relation between consecutive
diagonal matrices $\Lambda^{[m]}$, and the positivity and normalization of
$\Lambda^{[m]}$, are not formalized; only the existence statement with the bound
$D\le d^{\lfloor N/2\rfloor}$ is.

\section{Elimination plan}

The restrictions above concern the minimality of the printed size and the canonical
gauge of the completeness theorem. The latter is formalized by diagonalizing
the reduced density operators $\rho_m$ in the representation of item 2 and
rescaling the bond bases by the square roots of their eigenvalues.

\bibliographystyle{alpha}
\bibliography{references}

\end{document}
